\section{Functions and Procedures}

This section details the business logic automation implemented through MySQL stored functions, procedures, and triggers.

\subsection{Account Management Functions}

\subsubsection{Function: Calculate Customer Total Spent}

\paragraph{Purpose:} Calculates the total amount spent by a customer across all non-cancelled orders, including shipping fees and voucher discounts.

\begin{lstlisting}[language=SQL, caption=Function: fn\_get\_customer\_total\_spent, label=lst:fn_customer_spent]
CREATE FUNCTION fn_get_customer_total_spent(p_customer_id INT)
RETURNS DECIMAL(15, 2)
DETERMINISTIC
READS SQL DATA
BEGIN
    DECLARE v_total_spent DECIMAL(15, 2);
    
    SELECT COALESCE(SUM(fn_calculate_order_total(o.order_id)), 0)
    INTO v_total_spent
    FROM `Order` o
    WHERE o.customer_id = p_customer_id
      AND o.status != 'Cancelled';
    
    RETURN v_total_spent;
END$$
\end{lstlisting}

\paragraph{Related Tables:} ACCOUNT, CUSTOMER, ORDER, ORDERITEM, SHIPPING, VOUCHER

\paragraph{Explanation:}
This function sums all completed order totals for a specific customer. It queries the ORDER table filtered by customer\_id and status != 'Cancelled', then aggregates order totals including shipping fees and subtracting applied voucher discounts. The function is DETERMINISTIC (same input always returns same output) and reads SQL data only.

\paragraph{Example Usage:}

\begin{lstlisting}[language=SQL]
-- Get total spent by customer ID 1
SELECT fn_get_customer_total_spent(1);
-- Result: 2,069,500.00 VND
-- Customer 1 orders: Order 1 (1,985,000) + Order 4 (84,500)

-- Get total spent by customer ID 2
SELECT fn_get_customer_total_spent(2);
-- Result: 0.00 VND
-- Customer 2 only has cancelled orders

-- Get total spent by customer ID 3
SELECT fn_get_customer_total_spent(3);
-- Result: 31,000.00 VND
-- Customer 3: Order 3 (31,000)

-- Get customer info with total spent
SELECT c.customer_code, a.full_name, 
       fn_get_customer_total_spent(c.customer_id) AS total_spent
FROM Customer c
INNER JOIN Account a ON c.customer_id = a.account_id
WHERE c.customer_id = 1;
\end{lstlisting}

\subsection{Shop Management Functions}

\subsubsection{Function: Calculate Shop Revenue}

\paragraph{Purpose:} Calculates total revenue generated by a shop from all non-cancelled orders, accounting for multi-seller scenarios.

\begin{lstlisting}[language=SQL, caption=Function: fn\_get\_shop\_revenue, label=lst:fn_shop_revenue]
CREATE FUNCTION fn_get_shop_revenue(p_shop_id INT)
RETURNS DECIMAL(15, 2)
DETERMINISTIC
READS SQL DATA
BEGIN
    DECLARE v_revenue DECIMAL(15, 2);
    DECLARE v_subtotal DECIMAL(15, 2);
    DECLARE v_discount_total DECIMAL(15, 2);
    
    -- Calculate shop revenue (OrderItem + Shipping - Voucher discount)
    -- Only count non-cancelled orders
    
    -- Calculate subtotal from OrderItem
    SELECT COALESCE(SUM(oi.quantity * oi.price_at_purchase), 0)
    INTO v_subtotal
    FROM OrderItem oi
    INNER JOIN `Order` o ON oi.order_id = o.order_id
    WHERE oi.shop_id = p_shop_id AND o.status != 'Cancelled';
    
    -- Calculate voucher discount (only if min_order_value met)
    SELECT COALESCE(SUM(
        CASE 
            WHEN v.discount_type = 'Percentage' AND v_subtotal >= COALESCE(v.min_order_value, 0) 
                 THEN v_subtotal * v.discount_value / 100
            WHEN v.discount_type = 'Amount' AND v_subtotal >= COALESCE(v.min_order_value, 0)
                 THEN v.discount_value
            ELSE 0
        END
    ), 0)
    INTO v_discount_total
    FROM `Order` o
    LEFT JOIN Voucher v ON o.voucher_id = v.voucher_id
    WHERE o.order_id IN (
        SELECT DISTINCT oi.order_id 
        FROM OrderItem oi 
        WHERE oi.shop_id = p_shop_id
    )
    AND o.status != 'Cancelled';
    
    SET v_revenue = v_subtotal - v_discount_total;
    
    IF v_revenue < 0 THEN
        SET v_revenue = 0;
    END IF;
    
    RETURN v_revenue;
END$$
\end{lstlisting}

\paragraph{Related Tables:} SHOP, ORDERITEM, ORDER

\paragraph{Explanation:}
This function calculates shop revenue by summing all OrderItem line items (quantity × price\_at\_purchase) where the shop\_id matches and the order status is not 'Cancelled'. Since orders can contain items from multiple shops, using shop\_id in OrderItem ensures accurate shop-specific revenue calculation.

\paragraph{Example Usage:}

\begin{lstlisting}[language=SQL]
-- Get revenue for shop ID 4
SELECT fn_get_shop_revenue(4);
-- Result: 1,980,000.00 VND
-- Calculation: Shop 4 items (1,000,000 + 1,200,000) - voucher (220,000)

-- Get revenue for shop ID 5
SELECT fn_get_shop_revenue(5);
-- Result: 25,000.00 VND

-- Get revenue for shop ID 6
SELECT fn_get_shop_revenue(6);
-- Result: 80,000.00 VND

-- Top shops by revenue
SELECT s.shop_code, s.shop_name, fn_get_shop_revenue(s.shop_id) AS revenue
FROM Shop s
ORDER BY fn_get_shop_revenue(s.shop_id) DESC
LIMIT 5;
\end{lstlisting}

\subsubsection{Function: Calculate Shop Rating}

\paragraph{Purpose:} Calculates the average rating of a shop based on verified customer reviews.

\begin{lstlisting}[language=SQL, caption=Function: fn\_calculate\_shop\_rating, label=lst:fn_shop_rating]
CREATE FUNCTION fn_calculate_shop_rating(p_shop_id INT)
RETURNS DECIMAL(3, 2)
DETERMINISTIC
READS SQL DATA
BEGIN
    DECLARE v_avg_rating DECIMAL(3, 2);
    
    SELECT COALESCE(AVG(r.rating), 0)
    INTO v_avg_rating
    FROM Review r
    WHERE r.target_type = 'Shop'
      AND r.target_id = p_shop_id;
    
    RETURN v_avg_rating;
END$$
\end{lstlisting}

\paragraph{Related Tables:} SHOP, REVIEW

\paragraph{Explanation:}
This function queries the REVIEW table where target\_type='Shop' and target\_id matches the shop\_id, then calculates the average rating. If no reviews exist, it returns 0. The function is called by a trigger whenever a new review is added to keep the shop rating current.

\paragraph{Example Usage:}

\begin{lstlisting}[language=SQL]
-- Get shop rating for shop 4
SELECT fn_calculate_shop_rating(4);
-- Result: 4.00
-- Reviews for Shop 4: (5 + 3) / 2 = 4.00

-- Get shop rating for shop 5
SELECT fn_calculate_shop_rating(5);
-- Result: 0.00
-- No reviews for Shop 5

-- Get shop rating for shop 6
SELECT fn_calculate_shop_rating(6);
-- Result: 5.00
-- Review for Shop 6: 5 stars

-- Get shop info with calculated rating
SELECT s.shop_code, s.shop_name, s.rating AS stored_rating,
       fn_calculate_shop_rating(s.shop_id) AS calculated_rating
FROM Shop s
WHERE s.shop_id = 4;
\end{lstlisting}

---

\subsection{Product Management Functions}

\subsubsection{Function: Calculate Product Rating}

\paragraph{Purpose:} Calculates the average rating of a product based on verified customer reviews.

\begin{lstlisting}[language=SQL, caption=Function: fn\_calculate\_product\_rating, label=lst:fn_product_rating]
CREATE FUNCTION fn_calculate_product_rating(p_product_id INT)
RETURNS DECIMAL(3, 2)
DETERMINISTIC
READS SQL DATA
BEGIN
    DECLARE v_avg_rating DECIMAL(3, 2);
    
    SELECT COALESCE(AVG(r.rating), 0)
    INTO v_avg_rating
    FROM Review r
    WHERE r.target_type = 'Product'
      AND r.target_id = p_product_id;
    
    RETURN v_avg_rating;
END$$
\end{lstlisting}

\paragraph{Related Tables:} PRODUCT, REVIEW

\paragraph{Explanation:}
Similar to shop rating function, this calculates product average rating from REVIEW table where target\_type='Product'. Returns 0 if no reviews exist. Used by triggers to maintain denormalized rating in PRODUCT table for query performance.

\paragraph{Example Usage:}

\begin{lstlisting}[language=SQL]
-- Get rating for product 1
SELECT fn_calculate_product_rating(1);
-- Result: 4.50
-- Reviews for Product 1: (5 + 4) / 2 = 4.50

-- Get rating for product 3
SELECT fn_calculate_product_rating(3);
-- Result: 4.00
-- Review for Product 3: 4 stars

-- Get rating for product 4
SELECT fn_calculate_product_rating(4);
-- Result: 5.00
-- Review for Product 4: 5 stars

-- Get top-rated products
SELECT p.product_id, p.product_name, 
       fn_calculate_product_rating(p.product_id) AS rating,
       COUNT(r.review_id) AS review_count
FROM Product p
LEFT JOIN Review r ON p.product_id = r.target_id 
                  AND r.target_type = 'Product'
GROUP BY p.product_id
ORDER BY rating DESC
LIMIT 10;
\end{lstlisting}

\subsubsection{Function: Calculate Order Total}

\paragraph{Purpose:} Calculates the final order total including subtotal, shipping fee, and voucher discount with minimum order value validation.

\begin{lstlisting}[language=SQL, caption=Function: fn\_calculate\_order\_total, label=lst:fn_order_total]
CREATE FUNCTION fn_calculate_order_total(p_order_id INT)
RETURNS DECIMAL(15, 2)
DETERMINISTIC
READS SQL DATA
BEGIN
    DECLARE v_subtotal DECIMAL(15, 2);
    DECLARE v_shipping_fee DECIMAL(10, 2);
    DECLARE v_discount DECIMAL(10, 2);
    DECLARE v_discount_type VARCHAR(20);
    DECLARE v_discount_value DECIMAL(10, 2);
    DECLARE v_min_order_value DECIMAL(10, 2);
    DECLARE v_total DECIMAL(15, 2);
    
    -- Calculate subtotal from order items
    SELECT COALESCE(SUM(oi.quantity * oi.price_at_purchase), 0)
    INTO v_subtotal 
    FROM OrderItem oi
    WHERE oi.order_id = p_order_id;
    
    -- Get shipping fee
    SELECT COALESCE(sh.fee, 0)
    INTO v_shipping_fee
    FROM `Order` o
    LEFT JOIN Shipping sh ON o.shipping_id = sh.shipping_id
    WHERE o.order_id = p_order_id;
    
    -- Get voucher discount with min_order_value check
    SELECT v.discount_type, COALESCE(v.discount_value, 0), COALESCE(v.min_order_value, 0)
    INTO v_discount_type, v_discount_value, v_min_order_value
    FROM `Order` o
    LEFT JOIN Voucher v ON o.voucher_id = v.voucher_id
    WHERE o.order_id = p_order_id;
    
    -- Calculate discount amount (only if subtotal meets minimum requirement)
    SET v_discount = 0;
    IF v_discount_type IS NOT NULL AND v_subtotal >= v_min_order_value THEN
        IF v_discount_type = 'Percentage' THEN
            SET v_discount = v_subtotal * (v_discount_value / 100);
        ELSEIF v_discount_type = 'Amount' THEN
            SET v_discount = v_discount_value;
        END IF;
    END IF;
    
    -- Calculate total = subtotal + shipping - discount
    SET v_total = v_subtotal + v_shipping_fee - v_discount;
    
    RETURN v_total;
END$$
\end{lstlisting}

\paragraph{Related Tables:} ORDER, ORDERITEM, SHIPPING, VOUCHER

\paragraph{Explanation:}
This comprehensive function calculates order totals by:
\begin{enumerate}
    \item Summing all OrderItem line items (quantity × price\_at\_purchase)
    \item Adding shipping fee from related Shipping record
    \item Calculating voucher discount (percentage or fixed amount) with min\_order\_value validation
    \item Subtracting discount from subtotal + shipping
    \item Ensuring final total is non-negative
\end{enumerate}

\paragraph{Example Usage:}

\begin{lstlisting}[language=SQL]
-- Calculate order total for order 1
SELECT fn_calculate_order_total(1);
-- Result: 1,985,000.00 VND
-- Calculation: (1,000,000 + 1,200,000) + 5,000 - 220,000

-- Calculate order total for order 2
SELECT fn_calculate_order_total(2);
-- Result: 904,000.00 VND
-- Calculation: 900,000 + 4,000

-- Get order summary with calculated total
SELECT o.order_id, o.order_date, a.full_name,
       fn_calculate_order_total(o.order_id) AS total_amount,
       o.status, o.payment_method
FROM `Order` o
INNER JOIN Customer c ON o.customer_id = c.customer_id
INNER JOIN Account a ON c.customer_id = a.account_id
WHERE o.order_id = 1;
\end{lstlisting}

\subsubsection{Function: Get Order Total Items with Cursor}

\paragraph{Purpose:} Demonstrates cursor usage to loop through OrderItems and calculate total quantity in an order.

\begin{lstlisting}[language=SQL, caption=Function: fn\_get\_order\_total\_items\_with\_cursor, label=lst:fn_cursor_items]
CREATE FUNCTION fn_get_order_total_items_with_cursor(p_order_id INT)
RETURNS INT
DETERMINISTIC
READS SQL DATA
BEGIN
    DECLARE v_total_items INT DEFAULT 0;
    DECLARE v_done INT DEFAULT FALSE;
    DECLARE v_item_quantity INT;
    
    -- Declare cursor for OrderItems
    DECLARE order_items_cursor CURSOR FOR
        SELECT quantity
        FROM OrderItem
        WHERE order_id = p_order_id;
    
    -- Declare continue handler for cursor
    DECLARE CONTINUE HANDLER FOR NOT FOUND SET v_done = TRUE;
    
    -- Validate order exists
    IF NOT EXISTS (SELECT 1 FROM `Order` WHERE order_id = p_order_id) THEN
        RETURN 0;
    END IF;
    
    -- Open cursor
    OPEN order_items_cursor;
    
    -- Loop through each item
    item_loop: LOOP
        FETCH order_items_cursor INTO v_item_quantity;
        
        IF v_done THEN
            LEAVE item_loop;
        END IF;
        
        SET v_total_items = v_total_items + v_item_quantity;
    END LOOP item_loop;
    
    -- Close cursor
    CLOSE order_items_cursor;
    
    RETURN v_total_items;
END$$
DELIMITER ;
\end{lstlisting}

\paragraph{Related Tables:} ORDER, ORDERITEM

\paragraph{Explanation:}
This function demonstrates advanced cursor usage in MySQL stored functions. It:
\begin{enumerate}
    \item Declares a cursor to iterate through all OrderItems for a specific order
    \item Uses CONTINUE HANDLER to catch NOT FOUND exception when cursor reaches end
    \item Loops through each item, accumulating total quantity
    \item Returns cumulative total items count
\end{enumerate}

\paragraph{Example Usage:}

\begin{lstlisting}[language=SQL]
-- Get total items in order 1
SELECT fn_get_order_total_items_with_cursor(1);
-- Result: 2
-- Order 1 has 2 items (qty 1 + qty 1 = 2 total units)

-- Get total items in order 2
SELECT fn_get_order_total_items_with_cursor(2);
-- Result: 1
-- Order 2 has 1 item (qty 1)

-- Get total items for non-existent order
SELECT fn_get_order_total_items_with_cursor(999);
-- Result: 0
-- Order 999 does not exist

-- Order summary with total items
SELECT o.order_id, 
       COUNT(oi.order_item_id) AS line_items,
       fn_get_order_total_items_with_cursor(o.order_id) AS total_items
FROM `Order` o
LEFT JOIN OrderItem oi ON o.order_id = oi.order_id
WHERE o.order_id = 1
GROUP BY o.order_id;
\end{lstlisting}

\subsection{Category Hierarchy Procedures}

\subsubsection{Procedure: Get Category Hierarchy (All Levels)}

\paragraph{Purpose:} Displays complete hierarchical structure of all categories using recursive CTE (Common Table Expression).

\begin{lstlisting}[language=SQL, caption=Procedure: sp\_get\_category\_hierarchy, label=lst:sp_category_hierarchy]
CREATE PROCEDURE sp_get_category_hierarchy()
BEGIN
    WITH RECURSIVE category_tree AS (
        -- Anchor: Get all root categories (parent_category_id IS NULL)
        SELECT 
            category_id,
            category_name,
            parent_category_id,
            0 AS level,
            CAST(category_name AS CHAR(500)) AS path
        FROM Category
        WHERE parent_category_id IS NULL
        
        UNION ALL
        
        -- Recursive: Get all child categories
        SELECT 
            c.category_id,
            c.category_name,
            c.parent_category_id,
            ct.level + 1,
            CONCAT(ct.path, ' > ', c.category_name)
        FROM Category c
        INNER JOIN category_tree ct ON c.parent_category_id = ct.category_id
        WHERE ct.level < 10  -- Prevent infinite loop, limit to 10 levels
    )
    SELECT 
        category_id,
        category_name,
        parent_category_id,
        level,
        path,
        REPEAT('  ', level) AS indent
    FROM category_tree
    ORDER BY path;
END$$
\end{lstlisting}

\paragraph{Related Tables:} CATEGORY

\paragraph{Explanation:}
This procedure uses MySQL's recursive CTE (WITH RECURSIVE) to traverse the category hierarchy:
\begin{enumerate}
    \item \textbf{Anchor Query}: Selects all root categories (parent\_category\_id IS NULL) at level 0
    \item \textbf{Recursive Query}: Joins child categories to the tree, incrementing level
    \item \textbf{Path Building}: Constructs a breadcrumb path (e.g., "Electronics > Phones > Smartphones")
    \item \textbf{Recursion Limit}: Stops at level 10 to prevent infinite loops
\end{enumerate}

\paragraph{Example Usage:}

\begin{lstlisting}[language=SQL]
-- Get complete category hierarchy
CALL sp_get_category_hierarchy();

-- Result:
-- category_id | category_name | parent_id | level | path
-- 1           | Electronics   | NULL      | 0     | Electronics
-- 3           | Phones        | 1         | 1     | Electronics > Phones
-- 4           | Laptops       | 1         | 1     | Electronics > Laptops
-- 2           | Fashion       | NULL      | 0     | Fashion
-- 5           | Clothes       | 2         | 1     | Fashion > Clothes
\end{lstlisting}

\subsubsection{Procedure: Get All Children of a Category}

\paragraph{Purpose:} Retrieves all descendant categories (children, grandchildren, etc.) of a specified parent category.

\begin{lstlisting}[language=SQL, caption=Procedure: sp\_get\_category\_all\_children, label=lst:sp_category_children]
CREATE PROCEDURE sp_get_category_all_children(IN p_category_id INT)
BEGIN
    IF NOT EXISTS (SELECT 1 FROM Category WHERE category_id = p_category_id) THEN
        SIGNAL SQLSTATE '45000' 
        SET MESSAGE_TEXT = 'Category does not exist';
    END IF;
    
    WITH RECURSIVE children_tree AS (
        -- Anchor: Get the root category
        SELECT 
            category_id,
            category_name,
            parent_category_id,
            0 AS level
        FROM Category
        WHERE category_id = p_category_id
        
        UNION ALL
        
        -- Recursive: Get all child categories at next level
        SELECT 
            c.category_id,
            c.category_name,
            c.parent_category_id,
            ct.level + 1
        FROM Category c
        INNER JOIN children_tree ct ON c.parent_category_id = ct.category_id
        WHERE ct.level < 10
    )
    SELECT 
        category_id,
        category_name,
        parent_category_id,
        level,
        REPEAT('  ', level) AS indent
    FROM children_tree
    ORDER BY level, category_name;
END$$
\end{lstlisting}

\paragraph{Related Tables:} CATEGORY

\paragraph{Explanation:}
This procedure finds all descendants of a category:
\begin{enumerate}
    \item Validates that the input category exists (raises error if not)
    \item Anchor query retrieves the starting category at level 0
    \item Recursive query finds all immediate children, then their children, etc.
    \item Returns hierarchically structured results with indentation
\end{enumerate}

\paragraph{Example Usage:}

\begin{lstlisting}[language=SQL]
-- Get all children of Electronics (category_id = 1)
CALL sp_get_category_all_children(1);

-- Result:
-- category_id | category_name | parent_id | level
-- 1           | Electronics   | NULL      | 0
-- 3           | Phones        | 1         | 1
-- 4           | Laptops       | 1         | 1

-- Get all children of Fashion (category_id = 2)
CALL sp_get_category_all_children(2);

-- Result:
-- category_id | category_name | parent_id | level
-- 2           | Fashion       | NULL      | 0
-- 5           | Clothes       | 2         | 1
\end{lstlisting}

\subsubsection{Procedure: Get All Parents of a Category}

\paragraph{Purpose:} Retrieves all ancestor categories (parent, grandparent, etc.) of a specified category, showing the reverse hierarchy path.

\begin{lstlisting}[language=SQL, caption=Procedure: sp\_get\_category\_all\_parents, label=lst:sp_category_parents]
CREATE PROCEDURE sp_get_category_all_parents(IN p_category_id INT)
BEGIN
    IF NOT EXISTS (SELECT 1 FROM Category WHERE category_id = p_category_id) THEN
        SIGNAL SQLSTATE '45000' 
        SET MESSAGE_TEXT = 'Category does not exist';
    END IF;
    
    WITH RECURSIVE parents_tree AS (
        -- Anchor: Get the root category
        SELECT 
            category_id,
            category_name,
            parent_category_id,
            0 AS level
        FROM Category
        WHERE category_id = p_category_id
        
        UNION ALL
        
        -- Recursive: Get all parent categories at upper level
        SELECT 
            c.category_id,
            c.category_name,
            c.parent_category_id,
            pt.level + 1
        FROM Category c
        INNER JOIN parents_tree pt ON pt.parent_category_id = c.category_id
        WHERE pt.level < 10
    )
    SELECT 
        category_id,
        category_name,
        parent_category_id,
        level,
        REPEAT('  ', level) AS indent
    FROM parents_tree
    ORDER BY level DESC, category_name;
END$$
\end{lstlisting}

\paragraph{Related Tables:} CATEGORY

\paragraph{Explanation:}
This procedure traverses the hierarchy upward to find all ancestors:
\begin{enumerate}
    \item Validates category existence
    \item Anchor retrieves the specified category
    \item Recursive query follows parent\_category\_id chain upward
    \item Returns results ordered from leaf to root (DESC level)
\end{enumerate}

\paragraph{Example Usage:}

\begin{lstlisting}[language=SQL]
-- Get all parents of Phones (category_id = 3)
CALL sp_get_category_all_parents(3);

-- Result (showing breadcrumb trail):
-- category_id | category_name | parent_id | level
-- 3           | Phones        | 1         | 0
-- 1           | Electronics   | NULL      | 1
-- This shows: Phones < Electronics

-- Get all parents of Clothes (category_id = 5)
CALL sp_get_category_all_parents(5);

-- Result:
-- category_id | category_name | parent_id | level
-- 5           | Clothes       | 2         | 0
-- 2           | Fashion       | NULL      | 1
-- This shows: Clothes < Fashion
\end{lstlisting}


\subsection{Order Management Procedures}

\subsubsection{Procedure: Create Order from Cart}

\paragraph{Purpose:} Complex business logic procedure that creates an order from customer's shopping cart, validates inventory, applies vouchers, and updates related tables.

\begin{lstlisting}[language=SQL, caption=Procedure: sp\_create\_order\_from\_cart (Part 1), label=lst:sp_create_order_part1]
CREATE PROCEDURE sp_create_order_from_cart(
    IN p_customer_id INT,
    IN p_shipping_id INT,
    IN p_voucher_id INT,
    IN p_shipping_address VARCHAR(255),
    IN p_payment_method VARCHAR(50),
    IN p_note TEXT,
    OUT p_order_id INT
)
BEGIN
    DECLARE v_subtotal DECIMAL(15, 2) DEFAULT 0;
    DECLARE v_shipping_fee DECIMAL(10, 2) DEFAULT 0;
    DECLARE v_discount DECIMAL(10, 2) DEFAULT 0;
    DECLARE v_total_amount DECIMAL(15, 2) DEFAULT 0;
    DECLARE v_cart_id INT;
    DECLARE v_cart_count INT DEFAULT 0;
    DECLARE v_discount_type VARCHAR(20);
    DECLARE v_discount_value DECIMAL(10, 2);
    DECLARE v_min_order_value DECIMAL(10, 2);
    DECLARE v_usage_limit INT;
    DECLARE v_used_count INT;
    
    -- Validation checks
    IF NOT EXISTS (SELECT 1 FROM Customer WHERE customer_id = p_customer_id) THEN
        SIGNAL SQLSTATE '45000' 
        SET MESSAGE_TEXT = 'Customer does not exist';
    END IF;
    
    IF NOT EXISTS (SELECT 1 FROM Shipping WHERE shipping_id = p_shipping_id 
                   AND status = 'Active') THEN
        SIGNAL SQLSTATE '45000' 
        SET MESSAGE_TEXT = 'Invalid or inactive shipping method';
    END IF;
    
    -- Get and validate cart
    SELECT cart_id INTO v_cart_id
    FROM Cart
    WHERE customer_id = p_customer_id;
    
    IF v_cart_id IS NULL THEN
        SIGNAL SQLSTATE '45000' 
        SET MESSAGE_TEXT = 'Cart not found';
    END IF;
    
    SELECT COUNT(*) INTO v_cart_count
    FROM CartItem
    WHERE cart_id = v_cart_id;
    
    IF v_cart_count = 0 THEN
        SIGNAL SQLSTATE '45000' 
        SET MESSAGE_TEXT = 'Cart is empty';
    END IF;
    
    -- Calculate order amounts
    SELECT COALESCE(SUM(ci.quantity * pi.price), 0)
    INTO v_subtotal
    FROM CartItem ci
    INNER JOIN ProductItem pi ON ci.item_id = pi.item_id
    WHERE ci.cart_id = v_cart_id;
    
    SELECT fee INTO v_shipping_fee
    FROM Shipping
    WHERE shipping_id = p_shipping_id;
\end{lstlisting}

\begin{lstlisting}[language=SQL, caption=Procedure: sp\_create\_order\_from\_cart (Part 2), label=lst:sp_create_order_part2]
    
    -- Apply voucher if provided
    IF p_voucher_id IS NOT NULL THEN
        SELECT discount_type, discount_value, min_order_value, 
               usage_limit, used_count
        INTO v_discount_type, v_discount_value, v_min_order_value, 
             v_usage_limit, v_used_count
        FROM Voucher
        WHERE voucher_id = p_voucher_id
          AND status = 'Active'
          AND expired_date > CURDATE();
        
        IF v_discount_type IS NULL THEN
            SIGNAL SQLSTATE '45000' 
            SET MESSAGE_TEXT = 'Voucher is invalid or expired';
        END IF;
        
        IF v_subtotal < v_min_order_value THEN
            SIGNAL SQLSTATE '45000' 
            SET MESSAGE_TEXT = 'Order value does not meet voucher minimum';
        END IF;
        
        IF v_usage_limit > 0 AND v_used_count >= v_usage_limit THEN
            SIGNAL SQLSTATE '45000' 
            SET MESSAGE_TEXT = 'Voucher usage limit reached';
        END IF;
        
        -- Calculate discount
        IF v_discount_type = 'Percentage' THEN
            SET v_discount = v_subtotal * (v_discount_value / 100);
        ELSEIF v_discount_type = 'Amount' THEN
            SET v_discount = v_discount_value;
        END IF;
        
        -- Update voucher usage
        UPDATE Voucher
        SET used_count = used_count + 1
        WHERE voucher_id = p_voucher_id;
    END IF;
    
    -- Calculate total
    SET v_total_amount = v_subtotal + v_shipping_fee - v_discount;
    IF v_total_amount < 0 THEN
        SET v_total_amount = 0;
    END IF;
    
    -- Create order
    INSERT INTO `Order` (
        customer_id, shipping_id, voucher_id, status,
        shipping_address, total_amount, payment_method, note
    ) VALUES (
        p_customer_id, p_shipping_id, p_voucher_id, 'Processing',
        p_shipping_address, v_total_amount, p_payment_method, p_note
    );
    
    SET p_order_id = LAST_INSERT_ID();
    
    -- Move cart items to order items
    INSERT INTO OrderItem (order_id, item_id, shop_id, quantity, price_at_purchase)
    SELECT p_order_id, ci.item_id, pi.shop_id, ci.quantity, pi.price
    FROM CartItem ci
    INNER JOIN ProductItem pi ON ci.item_id = pi.item_id
    WHERE ci.cart_id = v_cart_id;
    
    -- Clear cart
    DELETE FROM CartItem WHERE cart_id = v_cart_id;
    
    -- Update customer stats
    UPDATE Customer
    SET total_order = total_order + 1,
        total_spent = total_spent + v_total_amount
    WHERE customer_id = p_customer_id;
END$$
\end{lstlisting}

\paragraph{Related Tables:} CUSTOMER, CART, CARTITEM, ORDER, ORDERITEM, SHIPPING, VOUCHER, PRODUCTITEM

\paragraph{Explanation:}
This is a complex transaction procedure that:
\begin{enumerate}
    \item Validates customer, shipping method, and cart existence
    \item Checks that cart contains items
    \item Calculates subtotal from cart items
    \item Validates and applies voucher (with min order value check)
    \item Creates ORDER record with final total
    \item Converts CARTITEM records to ORDERITEM records
    \item Clears the shopping cart
    \item Updates customer statistics (total\_order, total\_spent)
\end{enumerate}

\paragraph{Example Usage:}

\begin{lstlisting}[language=SQL]
-- Setup: Add items to Customer 1's cart
DELETE FROM CartItem WHERE cart_id = 1;
INSERT INTO CartItem (cart_id, item_id, quantity) VALUES (1, 1, 1);
INSERT INTO CartItem (cart_id, item_id, quantity) VALUES (1, 4, 2);

-- Verify cart contents before order creation
SELECT ci.cart_id, ci.item_id, ci.quantity, pi.price 
FROM CartItem ci
INNER JOIN ProductItem pi ON ci.item_id = pi.item_id
WHERE ci.cart_id = 1;
-- Result: 2 rows (item 1, item 4)

-- Create order from customer 1's cart
CALL sp_create_order_from_cart(
    1,                          -- customer_id
    1,                          -- shipping_id
    NULL,                       -- voucher_id (no voucher)
    '123 Updated Street',       -- shipping_address
    'Banking',                  -- payment_method
    'Test from cart',           -- note
    @new_order_id              -- OUT parameter
);

SELECT @new_order_id AS order_id;
-- Result: New order_id (e.g., 5)

-- Verify order was created
SELECT order_id, customer_id, status, payment_method 
FROM \`Order\` WHERE order_id = @new_order_id;
-- Result: status='Processing', payment_method='Banking'

-- Verify cart is empty after order creation
SELECT COUNT(*) AS remaining_items FROM CartItem WHERE cart_id = 1;
-- Result: 0 (cart cleared after order creation)
\end{lstlisting}

\subsubsection{Procedure: Apply Voucher}

\paragraph{Purpose:} Validates a voucher code and calculates discount amount with comprehensive business rule checking.

\begin{lstlisting}[language=SQL, caption=Procedure: sp\_apply\_voucher, label=lst:sp_apply_voucher]
CREATE PROCEDURE sp_apply_voucher(
    IN p_voucher_code VARCHAR(50),
    IN p_order_amount DECIMAL(15, 2),
    OUT p_discount_amount DECIMAL(10, 2),
    OUT p_is_valid TINYINT,
    OUT p_message VARCHAR(255)
)
BEGIN
    DECLARE v_voucher_id INT;
    DECLARE v_discount_type VARCHAR(20);
    DECLARE v_discount_value DECIMAL(10, 2);
    DECLARE v_min_order_value DECIMAL(10, 2);
    DECLARE v_usage_limit INT;
    DECLARE v_used_count INT;
    DECLARE v_status VARCHAR(20);
    
    SET p_is_valid = 0;
    SET p_discount_amount = 0;
    
    -- Get voucher by code
    SELECT 
        voucher_id, discount_type, discount_value, min_order_value,
        usage_limit, used_count, status
    INTO 
        v_voucher_id, v_discount_type, v_discount_value, v_min_order_value,
        v_usage_limit, v_used_count, v_status
    FROM Voucher
    WHERE code = p_voucher_code;
    
    -- Validation chain
    IF v_voucher_id IS NULL THEN
        SET p_message = 'Voucher code not found';
    ELSEIF v_status != 'Active' THEN
        SET p_message = 'Voucher is not active';
    ELSEIF CURDATE() > (SELECT expired_date FROM Voucher WHERE voucher_id = v_voucher_id) THEN
        SET p_message = 'Voucher has expired';
    ELSEIF v_usage_limit > 0 AND v_used_count >= v_usage_limit THEN
        SET p_message = 'Voucher usage limit reached';
    ELSEIF p_order_amount < v_min_order_value THEN
        SET p_message = CONCAT('Minimum order value is ', v_min_order_value);
    ELSE
        -- Calculate discount
        CASE v_discount_type
            WHEN 'Percentage' THEN
                SET p_discount_amount = p_order_amount * (v_discount_value / 100);
            WHEN 'Amount' THEN
                SET p_discount_amount = v_discount_value;
        END CASE;
        
        -- Cap discount at order amount
        IF p_discount_amount > p_order_amount THEN
            SET p_discount_amount = p_order_amount;
        END IF;
        
        SET p_is_valid = 1;
        SET p_message = 'Voucher applied successfully';
    END IF;
END$$
\end{lstlisting}

\paragraph{Related Tables:} VOUCHER

\paragraph{Explanation:}
This procedure validates vouchers through a series of checks:
\begin{enumerate}
    \item Checks if voucher code exists in database
    \item Verifies voucher is in 'Active' status
    \item Checks expiration date against current date
    \item Validates usage count hasn't exceeded limit
    \item Validates order amount meets minimum requirement
    \item Calculates discount (percentage or fixed amount)
    \item Caps discount to not exceed order amount
\end{enumerate}

\paragraph{Example Usage:}

\begin{lstlisting}[language=SQL]
-- Test 1: Valid percentage voucher (DISCOUNT10)
SET @discount = 0;
SET @valid = FALSE;
SET @msg = '';
CALL sp_apply_voucher('DISCOUNT10', 500000, @discount, @valid, @msg);
SELECT @discount AS discount_amount, @valid AS is_valid, @msg AS message;
-- Result: discount=50000, is_valid=1, message='Voucher applied successfully'
-- DISCOUNT10 is 10% percentage discount

-- Test 2: Valid amount voucher (SAVE50) - Order meets minimum
CALL sp_apply_voucher('SAVE50', 300000, @discount, @valid, @msg);
SELECT @discount, @valid, @msg;
-- Result: discount=50000, is_valid=1, message='Voucher applied successfully'

-- Test 3: Order below minimum requirement
CALL sp_apply_voucher('SAVE50', 150000, @discount, @valid, @msg);
SELECT @msg;
-- Result: message mentions minimum order value not met

-- Test 4: Non-existent voucher
CALL sp_apply_voucher('NONEXISTENT', 500000, @discount, @valid, @msg);
SELECT @msg;
-- Result: 'Voucher code not found'
\end{lstlisting}

\subsection{Shop Management Procedures}

\subsubsection{Procedure: Get Shop Revenue Report}

\paragraph{Purpose:} Comprehensive reporting procedure that generates shop revenue summary and top-selling products for a date range.

\begin{lstlisting}[language=SQL, caption=Procedure: sp\_get\_shop\_revenue\_report, label=lst:sp_shop_report_part1]
CREATE PROCEDURE sp_get_shop_revenue_report(
    IN p_shop_id INT,
    IN p_from_date DATE,
    IN p_to_date DATE
)
BEGIN
    -- Validate shop exists
    IF NOT EXISTS (SELECT 1 FROM Shop WHERE shop_id = p_shop_id) THEN
        SIGNAL SQLSTATE '45000' 
        SET MESSAGE_TEXT = 'Shop does not exist';
    END IF;
    
    -- Main revenue summary
    SELECT 
        s.shop_id,
        s.shop_code,
        s.shop_name,
        COUNT(DISTINCT o.order_id) AS total_orders,
        COUNT(DISTINCT oi.order_item_id) AS total_line_items,
        SUM(oi.quantity) AS total_units_sold,
        COALESCE(SUM(oi.quantity * oi.price_at_purchase), 0) AS total_revenue,
        COALESCE(AVG(oi.quantity * oi.price_at_purchase), 0) AS avg_order_value,
        s.rating AS shop_rating,
        p_from_date AS report_from,
        p_to_date AS report_to
    FROM Shop s
    LEFT JOIN OrderItem oi ON s.shop_id = oi.shop_id
    LEFT JOIN `Order` o ON oi.order_id = o.order_id
        AND o.order_date BETWEEN p_from_date AND DATE_ADD(p_to_date, INTERVAL 1 DAY)
        AND o.status != 'Cancelled'
    WHERE s.shop_id = p_shop_id
    GROUP BY s.shop_id, s.shop_code, s.shop_name, s.rating;
    
    -- Top 10 selling product items
    SELECT 
        p.product_id,
        p.product_name,
        pi.color,
        pi.type,
        COUNT(oi.order_item_id) AS times_ordered,
        SUM(oi.quantity) AS total_quantity_sold,
        SUM(oi.quantity * oi.price_at_purchase) AS item_revenue,
        ROUND(SUM(oi.quantity * oi.price_at_purchase) / COUNT(oi.order_item_id), 2) 
            AS avg_item_value
    FROM Product p
    INNER JOIN ProductItem pi ON p.product_id = pi.product_id
    INNER JOIN OrderItem oi ON pi.item_id = oi.item_id
    INNER JOIN `Order` o ON oi.order_id = o.order_id
        AND o.order_date BETWEEN p_from_date AND DATE_ADD(p_to_date, INTERVAL 1 DAY)
        AND o.status != 'Cancelled'
    WHERE pi.shop_id = p_shop_id
    GROUP BY p.product_id, p.product_name, pi.color, pi.type
    ORDER BY item_revenue DESC
    LIMIT 10;
END$$
\end{lstlisting}

\paragraph{Related Tables:} SHOP, ORDER, ORDERITEM, PRODUCT, PRODUCTITEM

\paragraph{Explanation:}
This reporting procedure generates two result sets:

\paragraph{Result Set 1 - Shop Summary:}
\begin{itemize}
    \item Total orders during period
    \item Line items and units sold
    \item Revenue metrics (total and average)
    \item Current shop rating
    \item Date range of report
\end{itemize}

\paragraph{Result Set 2 - Top 10 Products:}
\begin{itemize}
    \item Product and variant details (color, type)
    \item Order frequency
    \item Revenue contribution
    \item Average item value
\end{itemize}

\paragraph{Example Usage:}

\begin{lstlisting}[language=SQL]
-- Get revenue report for shop 4 for full year 2025
CALL sp_get_shop_revenue_report(4, '2025-01-01', '2025-12-31');

-- Result Set 1 (Summary):
-- shop_id | shop_code | shop_name | total_orders | total_units | revenue
-- 4       | SHOP00004 | Shop ABC  | 1            | 2           | 2,200,000

-- Result Set 2 (Top Products):
-- product_name | color | type  | times_ordered | qty_sold | revenue
-- iPhone 15    | Red   | 128GB | 1             | 1        | 1,000,000
-- iPhone 15    | Blue  | 256GB | 1             | 1        | 1,200,000

-- Get revenue report for shop 5
CALL sp_get_shop_revenue_report(5, '2025-01-01', '2025-12-31');

-- Result Set 1:
-- shop_id | shop_code | shop_name | total_orders | total_units | revenue
-- 5       | SHOP00005 | Shop XYZ  | 1            | 1           | 25,000

-- Result Set 2:
-- product_name | color | type    | times_ordered | qty_sold | revenue
-- T-Shirt      | White | Size M  | 1             | 1        | 25,000
\end{lstlisting}

\subsection{Product CRUD Procedures}

\subsubsection{Procedure: Add Product}

\paragraph{Purpose:} Creates a new product in the system with validation for shop and category existence.

\begin{lstlisting}[language=SQL, caption=Procedure: sp\_add\_product, label=lst:sp_add_product]
CREATE PROCEDURE sp_add_product(
    IN p_shop_id INT,
    IN p_category_id INT,
    IN p_product_name VARCHAR(255),
    IN p_description TEXT,
    IN p_image VARCHAR(255),
    OUT p_product_id INT,
    OUT p_status VARCHAR(50)
)
BEGIN
    DECLARE EXIT HANDLER FOR SQLEXCEPTION
    BEGIN
        SET p_product_id = -1;
        SET p_status = 'Error: Invalid shop or category ID';
    END;
    
    SET p_product_id = -1;
    SET p_status = 'Error';
    
    -- Validate shop exists
    IF NOT EXISTS (SELECT 1 FROM Shop WHERE shop_id = p_shop_id) THEN
        SET p_status = 'Error: Shop not found';
    ELSEIF NOT EXISTS (SELECT 1 FROM Category 
                       WHERE category_id = p_category_id) THEN
        SET p_status = 'Error: Category not found';
    ELSE
        -- Insert product
        INSERT INTO Product (shop_id, category_id, product_name, 
                            description, image, status)
        VALUES (p_shop_id, p_category_id, p_product_name, 
                p_description, p_image, 'In stock');
        
        SET p_product_id = LAST_INSERT_ID();
        SET p_status = 'Success';
    END IF;
END$$
\end{lstlisting}

\paragraph{Related Tables:} PRODUCT, SHOP, CATEGORY

\paragraph{Explanation:}
This procedure creates a new product with error handling:
\begin{enumerate}
    \item \textbf{EXIT HANDLER}: Catches SQL exceptions and returns error status
    \item \textbf{Shop Validation}: Verifies shop\_id exists in Shop table
    \item \textbf{Category Validation}: Verifies category\_id exists in Category table
    \item \textbf{Product Creation}: Inserts new product with default status 'In stock'
    \item \textbf{OUT Parameters}: Returns new product\_id and status message
\end{enumerate}

\paragraph{Example Usage:}

\begin{lstlisting}[language=SQL]
-- Add new product to Shop 4, Category 3 (Phones)
CALL sp_add_product(
    4,                              -- shop_id
    3,                              -- category_id
    'Samsung Galaxy S24',           -- product_name
    'Latest Samsung flagship',      -- description
    'samsung-s24.jpg',             -- image
    @product_id,                   -- OUT: product_id
    @status                        -- OUT: status message
);

SELECT @product_id, @status;
-- Result: @product_id = 5, @status = 'Success'

-- Verify product was created
SELECT product_id, product_name, status 
FROM Product WHERE product_id = @product_id;
-- Result: Samsung Galaxy S24, In stock

-- Test invalid shop ID
CALL sp_add_product(
    999,                           -- invalid shop_id
    3,
    'Test Product',
    'Description',
    'test.jpg',
    @product_id,
    @status
);

SELECT @status;
-- Result: 'Error: Shop not found'
\end{lstlisting}

\subsubsection{Procedure: Update Product}

\paragraph{Purpose:} Updates existing product information with NULL-safe field updates.

\begin{lstlisting}[language=SQL, caption=Procedure: sp\_update\_product, label=lst:sp_update_product]
CREATE PROCEDURE sp_update_product(
    IN p_product_id INT,
    IN p_product_name VARCHAR(255),
    IN p_description TEXT,
    IN p_image VARCHAR(255),
    IN p_status VARCHAR(50),
    OUT p_success BOOLEAN,
    OUT p_message VARCHAR(100)
)
BEGIN
    DECLARE EXIT HANDLER FOR SQLEXCEPTION
    BEGIN
        SET p_success = FALSE;
        SET p_message = 'Error: Database error';
    END;
    
    IF NOT EXISTS (SELECT 1 FROM Product WHERE product_id = p_product_id) THEN
        SET p_success = FALSE;
        SET p_message = 'Error: Product not found';
    ELSE
        UPDATE Product
        SET product_name = COALESCE(p_product_name, product_name),
            description = COALESCE(p_description, description),
            image = COALESCE(p_image, image),
            status = COALESCE(p_status, status)
        WHERE product_id = p_product_id;
        
        SET p_success = TRUE;
        SET p_message = 'Product updated successfully';
    END IF;
END$$
\end{lstlisting}

\paragraph{Related Tables:} PRODUCT

\paragraph{Explanation:}
This procedure demonstrates selective field updates:
\begin{enumerate}
    \item \textbf{COALESCE Pattern}: Updates only non-NULL input fields
    \item \textbf{Existence Check}: Validates product\_id before update
    \item \textbf{Error Handling}: Returns success flag and descriptive message
    \item \textbf{Partial Updates}: Allows updating individual fields without affecting others
\end{enumerate}

\paragraph{Example Usage:}

\begin{lstlisting}[language=SQL]
-- Update only product name and status
CALL sp_update_product(
    1,                              -- product_id
    'iPhone 15 Pro Max',           -- new name
    NULL,                          -- keep existing description
    NULL,                          -- keep existing image
    'Out of stock',                -- new status
    @success,
    @message
);

SELECT @success, @message;
-- Result: TRUE, 'Product updated successfully'

-- Verify update
SELECT product_name, status FROM Product WHERE product_id = 1;
-- Result: iPhone 15 Pro Max, Out of stock

-- Update only description
CALL sp_update_product(
    1,
    NULL,                          -- keep existing name
    'Updated description text',    -- new description
    NULL,
    NULL,
    @success,
    @message
);
\end{lstlisting}

\subsubsection{Procedure: Delete Product}

\paragraph{Purpose:} Deletes a product and all related product variants with CASCADE behavior.

\begin{lstlisting}[language=SQL, caption=Procedure: sp\_delete\_product, label=lst:sp_delete_product]
CREATE PROCEDURE sp_delete_product(
    IN p_product_id INT,
    OUT p_success BOOLEAN,
    OUT p_message VARCHAR(100)
)
BEGIN
    DECLARE v_affected_items INT;
    
    DECLARE EXIT HANDLER FOR SQLEXCEPTION
    BEGIN
        SET p_success = FALSE;
        SET p_message = 'Error: Database error';
    END;
    
    IF NOT EXISTS (SELECT 1 FROM Product WHERE product_id = p_product_id) THEN
        SET p_success = FALSE;
        SET p_message = 'Error: Product not found';
    ELSE
        -- Get count of affected items before deletion
        SELECT COUNT(*) INTO v_affected_items 
        FROM ProductItem 
        WHERE product_id = p_product_id;
        
        -- Delete product (CASCADE deletes ProductItems)
        DELETE FROM Product WHERE product_id = p_product_id;
        
        SET p_success = TRUE;
        SET p_message = CONCAT('Product deleted. Removed ', 
                               v_affected_items, ' variant(s)');
    END IF;
END$$
\end{lstlisting}

\paragraph{Related Tables:} PRODUCT, PRODUCTITEM

\paragraph{Explanation:}
This procedure handles cascading deletion:
\begin{enumerate}
    \item \textbf{Pre-deletion Count}: Records number of variants before deletion
    \item \textbf{CASCADE Behavior}: Foreign key constraints automatically delete related ProductItems
    \item \textbf{Informative Message}: Returns count of deleted variants
    \item \textbf{Transactional Safety}: EXIT HANDLER ensures atomicity
\end{enumerate}

\paragraph{Example Usage:}

\begin{lstlisting}[language=SQL]
-- Check product variants before deletion
SELECT COUNT(*) AS variant_count 
FROM ProductItem 
WHERE product_id = 1;
-- Result: 2 variants (Red 128GB, Blue 256GB)

-- Delete product
CALL sp_delete_product(1, @success, @message);

SELECT @success, @message;
-- Result: TRUE, 'Product deleted. Removed 2 variant(s)'

-- Verify deletion
SELECT COUNT(*) FROM Product WHERE product_id = 1;
-- Result: 0

SELECT COUNT(*) FROM ProductItem WHERE product_id = 1;
-- Result: 0 (cascaded deletion)
\end{lstlisting}

\subsubsection{Procedure: Add Product Item (Variant)}

\paragraph{Purpose:} Adds a new product variant with color, type, price, and stock information.

\begin{lstlisting}[language=SQL, caption=Procedure: sp\_add\_product\_item, label=lst:sp_add_product_item]
CREATE PROCEDURE sp_add_product_item(
    IN p_product_id INT,
    IN p_shop_id INT,
    IN p_color VARCHAR(50),
    IN p_type VARCHAR(50),
    IN p_price DECIMAL(15,2),
    IN p_stock INT,
    IN p_image_url VARCHAR(255),
    OUT p_item_id INT,
    OUT p_status VARCHAR(50)
)
BEGIN
    DECLARE EXIT HANDLER FOR SQLEXCEPTION
    BEGIN
        SET p_item_id = -1;
        SET p_status = 'Error: Invalid input';
    END;
    
    SET p_item_id = -1;
    SET p_status = 'Error';
    
    IF NOT EXISTS (SELECT 1 FROM Product WHERE product_id = p_product_id) THEN
        SET p_status = 'Error: Product not found';
    ELSEIF p_price <= 0 THEN
        SET p_status = 'Error: Price must be > 0';
    ELSEIF p_stock < 0 THEN
        SET p_status = 'Error: Stock cannot be negative';
    ELSE
        INSERT INTO ProductItem (product_id, shop_id, color, type, 
                                 price, stock, image_url)
        VALUES (p_product_id, p_shop_id, p_color, p_type, 
                p_price, p_stock, p_image_url);
        
        SET p_item_id = LAST_INSERT_ID();
        SET p_status = 'Success';
    END IF;
END$$
\end{lstlisting}

\paragraph{Related Tables:} PRODUCTITEM, PRODUCT, SHOP

\paragraph{Explanation:}
This procedure creates product variants with validation:
\begin{enumerate}
    \item \textbf{Product Existence}: Validates parent product exists
    \item \textbf{Price Validation}: Ensures price > 0
    \item \textbf{Stock Validation}: Ensures stock >= 0
    \item \textbf{Variant Details}: Stores color, type (size/capacity), image
\end{enumerate}

\paragraph{Example Usage:}

\begin{lstlisting}[language=SQL]
-- Add new variant: iPhone 15 Green 512GB
CALL sp_add_product_item(
    1,                          -- product_id (iPhone 15)
    4,                          -- shop_id
    'Green',                    -- color
    '512GB',                    -- type
    1400000.00,                 -- price
    15,                         -- stock
    'iphone15-green-512.jpg',  -- image_url
    @item_id,
    @status
);

SELECT @item_id, @status;
-- Result: @item_id = 5, @status = 'Success'

-- Verify variant was created
SELECT item_id, color, type, price, stock 
FROM ProductItem WHERE item_id = @item_id;
-- Result: 5, Green, 512GB, 1400000.00, 15

-- Test invalid price
CALL sp_add_product_item(
    1, 4, 'Black', '128GB',
    -100.00,                    -- invalid price
    10, 'test.jpg',
    @item_id, @status
);

SELECT @status;
-- Result: 'Error: Price must be > 0'
\end{lstlisting}

\subsubsection{Procedure: Update Product Item}

\paragraph{Purpose:} Updates product variant price, stock, and image with validation.

\begin{lstlisting}[language=SQL, caption=Procedure: sp\_update\_product\_item, label=lst:sp_update_product_item]
CREATE PROCEDURE sp_update_product_item(
    IN p_item_id INT,
    IN p_price DECIMAL(15,2),
    IN p_stock INT,
    IN p_image_url VARCHAR(255),
    OUT p_success BOOLEAN,
    OUT p_message VARCHAR(100)
)
BEGIN
    DECLARE EXIT HANDLER FOR SQLEXCEPTION
    BEGIN
        SET p_success = FALSE;
        SET p_message = 'Error: Database error';
    END;
    
    IF NOT EXISTS (SELECT 1 FROM ProductItem WHERE item_id = p_item_id) THEN
        SET p_success = FALSE;
        SET p_message = 'Error: Item not found';
    ELSEIF p_price IS NOT NULL AND p_price <= 0 THEN
        SET p_success = FALSE;
        SET p_message = 'Error: Price must be > 0';
    ELSEIF p_stock IS NOT NULL AND p_stock < 0 THEN
        SET p_success = FALSE;
        SET p_message = 'Error: Stock cannot be negative';
    ELSE
        UPDATE ProductItem
        SET price = COALESCE(p_price, price),
            stock = COALESCE(p_stock, stock),
            image_url = COALESCE(p_image_url, image_url)
        WHERE item_id = p_item_id;
        
        SET p_success = TRUE;
        SET p_message = 'Item updated successfully';
    END IF;
END$$
\end{lstlisting}

\paragraph{Related Tables:} PRODUCTITEM

\paragraph{Explanation:}
This procedure demonstrates selective updates with validation:
\begin{enumerate}
    \item \textbf{COALESCE Updates}: Updates only provided fields
    \item \textbf{Pre-update Validation}: Checks constraints before UPDATE
    \item \textbf{Business Rules}: Price must be positive, stock non-negative
\end{enumerate}

\paragraph{Example Usage:}

\begin{lstlisting}[language=SQL]
-- Update only price and stock
CALL sp_update_product_item(
    1,                          -- item_id
    1100000.00,                 -- new price
    50,                         -- new stock
    NULL,                       -- keep existing image
    @success,
    @message
);

SELECT @success, @message;
-- Result: TRUE, 'Item updated successfully'

-- Update only stock
CALL sp_update_product_item(
    1,
    NULL,                       -- keep existing price
    45,                         -- decrease stock
    NULL,
    @success,
    @message
);
\end{lstlisting}

\subsubsection{Procedure: Delete Product Item}

\paragraph{Purpose:} Deletes a specific product variant.

\begin{lstlisting}[language=SQL, caption=Procedure: sp\_delete\_product\_item, label=lst:sp_delete_product_item]
CREATE PROCEDURE sp_delete_product_item(
    IN p_item_id INT,
    OUT p_success BOOLEAN,
    OUT p_message VARCHAR(100)
)
BEGIN
    DECLARE EXIT HANDLER FOR SQLEXCEPTION
    BEGIN
        SET p_success = FALSE;
        SET p_message = 'Error: Database error';
    END;
    
    IF NOT EXISTS (SELECT 1 FROM ProductItem WHERE item_id = p_item_id) THEN
        SET p_success = FALSE;
        SET p_message = 'Error: Item not found';
    ELSE
        DELETE FROM ProductItem WHERE item_id = p_item_id;
        
        SET p_success = TRUE;
        SET p_message = 'Item deleted successfully';
    END IF;
END$$
\end{lstlisting}

\paragraph{Related Tables:} PRODUCTITEM

\paragraph{Explanation:}
Simple deletion procedure with existence validation and error handling.

\paragraph{Example Usage:}

\begin{lstlisting}[language=SQL]
-- Delete variant
CALL sp_delete_product_item(2, @success, @message);

SELECT @success, @message;
-- Result: TRUE, 'Item deleted successfully'

-- Verify deletion
SELECT COUNT(*) FROM ProductItem WHERE item_id = 2;
-- Result: 0
\end{lstlisting}

\subsection{Product Search Procedure}

\subsubsection{Procedure: Get Product List with Filters}

\paragraph{Purpose:} Dynamic product search with filters, sorting, aggregations, and pagination using dynamic SQL.

\begin{lstlisting}[language=SQL, caption=Procedure: sp\_get\_product\_list (Excerpt), label=lst:sp_product_list]
CREATE PROCEDURE sp_get_product_list(
    IN p_search VARCHAR(255),
    IN p_category_id INT,
    IN p_shop_id INT,
    IN p_min_price DECIMAL(15,2),
    IN p_max_price DECIMAL(15,2),
    IN p_status VARCHAR(50),
    IN p_sort_by VARCHAR(50),
    IN p_sort_order VARCHAR(10),
    IN p_page INT,
    IN p_limit INT,
    OUT p_total_count INT
)
BEGIN
    DECLARE v_offset INT;
    DECLARE v_order_clause VARCHAR(100);
    
    SET v_offset = (p_page - 1) * p_limit;
    
    -- Build ORDER clause
    SET v_order_clause = CASE p_sort_by
        WHEN 'price' THEN CONCAT('min_price ', COALESCE(p_sort_order, 'ASC'))
        WHEN 'rating' THEN CONCAT('avg_rating ', COALESCE(p_sort_order, 'DESC'))
        WHEN 'created_at' THEN CONCAT('p.created_at ', COALESCE(p_sort_order, 'DESC'))
        ELSE 'p.created_at DESC'
    END;
    
    -- Count total matching records
    SELECT COUNT(DISTINCT p.product_id) INTO p_total_count
    FROM Product p
    INNER JOIN Shop s ON p.shop_id = s.shop_id
    INNER JOIN Category c ON p.category_id = c.category_id
    LEFT JOIN ProductItem pi ON p.product_id = pi.product_id
    WHERE 1=1
        AND (p_search IS NULL OR p.product_name LIKE CONCAT('%', p_search, '%') 
             OR p.description LIKE CONCAT('%', p_search, '%'))
        AND (p_category_id IS NULL OR p.category_id = p_category_id)
        AND (p_shop_id IS NULL OR p.shop_id = p_shop_id)
        AND (p_status IS NULL OR p.status = p_status);
    
    -- Dynamic query construction
    SET @sql = CONCAT('
        SELECT 
            p.product_id,
            p.product_name,
            p.description,
            p.image,
            p.status,
            p.created_at,
            p.shop_id,
            s.shop_name,
            s.rating as shop_rating,
            c.category_id,
            c.category_name,
            MIN(pi.price) as min_price,
            MAX(pi.price) as max_price,
            SUM(pi.stock) as total_stock,
            COUNT(DISTINCT pi.item_id) as variant_count,
            COALESCE((
                SELECT AVG(r.rating)
                FROM Review r
                WHERE r.target_id = p.product_id AND r.target_type = "Product"
            ), 0) as avg_rating,
            COALESCE((
                SELECT COUNT(*)
                FROM Review r
                WHERE r.target_id = p.product_id AND r.target_type = "Product"
            ), 0) as review_count
        FROM Product p
        INNER JOIN Shop s ON p.shop_id = s.shop_id
        INNER JOIN Category c ON p.category_id = c.category_id
        LEFT JOIN ProductItem pi ON p.product_id = pi.product_id
        WHERE 1=1',
        CASE WHEN p_search IS NOT NULL THEN 
            CONCAT(' AND (p.product_name LIKE "%', p_search, '%" OR p.description LIKE "%', p_search, '%")')
        ELSE '' END,
        CASE WHEN p_category_id IS NOT NULL THEN CONCAT(' AND p.category_id = ', p_category_id) 
        ELSE '' END,
        CASE WHEN p_shop_id IS NOT NULL THEN CONCAT(' AND p.shop_id = ', p_shop_id) 
        ELSE '' END,
        CASE WHEN p_status IS NOT NULL THEN CONCAT(' AND p.status = "', p_status, '"') 
        ELSE '' END,
        ' GROUP BY p.product_id, p.product_name, p.description, p.image, p.status, 
                   p.created_at, p.shop_id, s.shop_name, s.rating, c.category_id, c.category_name',
        CASE WHEN p_min_price IS NOT NULL THEN CONCAT(' HAVING min_price >= ', p_min_price) 
        ELSE '' END,
        CASE WHEN p_max_price IS NOT NULL THEN 
            CONCAT(CASE WHEN p_min_price IS NOT NULL THEN ' AND' ELSE ' HAVING' END, 
                   ' max_price <= ', p_max_price)
        ELSE '' END,
        ' ORDER BY ', v_order_clause,
        ' LIMIT ', p_limit, ' OFFSET ', v_offset
    );
    
    PREPARE stmt FROM @sql;
    EXECUTE stmt;
    DEALLOCATE PREPARE stmt;
END$$
\end{lstlisting}

\paragraph{Related Tables:} PRODUCT, PRODUCTITEM, SHOP, CATEGORY, REVIEW

\paragraph{Explanation:}
This advanced procedure demonstrates:
\begin{enumerate}
    \item \textbf{Dynamic SQL Construction}: Builds query based on provided filters
    \item \textbf{Multiple Filters}: Search text, category, shop, status, price range
    \item \textbf{Aggregations}: MIN/MAX price, SUM stock, COUNT variants, AVG rating
    \item \textbf{Subqueries}: Rating and review count calculated via subselects
    \item \textbf{Sorting}: Dynamic sort order based on price/rating/date
    \item \textbf{Pagination}: Offset and limit for result set paging
    \item \textbf{Total Count}: OUT parameter returns total matching records
\end{enumerate}

\paragraph{Example Usage:}

\begin{lstlisting}[language=SQL]
-- Search smartphones under 10M, sorted by rating, page 1
DECLARE @total INT;
CALL sp_get_product_list(
    'iPhone',       -- search term
    3,              -- category_id (Smartphones)
    NULL,           -- shop_id (any)
    0,              -- min_price
    10000000,       -- max_price
    'In stock',     -- status
    'rating',       -- sort by rating
    'DESC',         -- descending
    1,              -- page 1
    20,             -- 20 items per page
    @total          -- OUT: total count
);

SELECT @total;  -- e.g., 45 total matches
\end{lstlisting}
